\section{两个泥潭：长期漂浮与算不过账}

上一章讲软硬件和数据，本章进入创业公司的风险。王鹤认为，具身智能公司如果陷入两个泥潭，天花板会很有限。第一个是“长期漂浮”：公司长时间停留在概念、演示和融资叙事里，不进入真实场景，不接受客户、成本和可靠性检验。第二个是“算不过来账”：即使进入场景，边际成本不下降，越交付越重，最终生产力不能成立。

这两个泥潭本质上都是反馈错位。长期漂浮意味着公司的 reward model 是 attention 和 valuation；算不过账意味着公司的技术路线没有让单位交付成本下降。具身智能是物理行业，不能只看模型 benchmark，也不能只看 demo 视频。它必须回答：一台机器人进入一个场景后，是否能比人更稳定、更便宜或更可扩展地完成一类任务？

\begin{figure}[H]
\centering
\includegraphics[width=0.94\textwidth]{two-traps.png}
\caption{两个泥潭：长期漂浮与算不过账都会限制公司天花板。自制概念图，依据 01:44:34--01:55:17 对谈内容与视频描述整理。}
\end{figure}

\begin{knowledgebox}{读图：一个是叙事风险，一个是经济性风险}
长期漂浮的问题是脱离真实客户，算不过账的问题是交付越多越亏。读图时要把两边连起来看：只做 demo 会避免短期成本暴露，但也错过真实数据；进入场景后如果边际成本不降，数据飞轮也无法支撑商业闭环。
\end{knowledgebox}

\subsection{应用场景内的泛化}

本节解释银河通用为什么选择“应用场景内泛化”。通用机器人很诱人，但过早追求全场景泛化会让任务、硬件、数据和成本同时爆炸。更现实的路线是在一个具体应用场景里做足泛化：例如货架场景中不同商品、摆放、遮挡、缺货、补货和异常情况。它不是单点自动化，而是在有限世界内构建可扩展能力。

这个思路和大模型产品化也相似：先把任务域、反馈、成本和数据闭环跑通，再逐步扩大能力边界。对机器人来说，场景内泛化还有一个额外好处：一旦出货，真实数据会回流，系统可以从真实失败中学习。没有出货量，数据飞轮只是幻觉。

\begin{importantbox}{场景内泛化}
具身智能早期的有效路线可能不是“先做通用机器人”，而是在一个真实付费场景内，把物体、任务、异常和成本做出足够泛化。场景边界越清晰，数据闭环和商业闭环越容易被验证。
\end{importantbox}

\subsection{本章小结}

具身智能公司要避开两个泥潭：不要长期漂浮在叙事里，也不要进入场景后算不过账。真正的路线是找到可付费场景，在场景内泛化，并让出货回流数据。

